\documentclass[10pt,a4paper]{amsart}
\usepackage[margin=19mm]{geometry}
\usepackage{amsmath,amssymb,mathtools}
\usepackage[scr=rsfs,cal=euler]{mathalfa}
\usepackage{xfrac}
\usepackage[unicode,hidelinks,bookmarks=false]{hyperref}
\newcommand{\ie}{i.e.\ }
\newcommand{\cf}{cf.\ }
\newcommand{\resp}{respectively\ }
\newcommand{\Subsection}{\subsection}
\providecommand{\nobreakdash}{\nobreak\hbox{-}}
\setlength{\emergencystretch}{2em}
\multlinegap0pt
\usepackage{mathtools}
\usepackage{enumitem}
\hypersetup{
    colorlinks=true,
    linkcolor=blue,
    citecolor=magenta,
    urlcolor=blue
}
\numberwithin{equation}{section}
\newtheorem{theorem}{Theorem}[section]
\newtheorem{proposition}[theorem]{Proposition}
\newtheorem{lemma}[theorem]{Lemma}
\newtheorem{corollary}[theorem]{Corollary}
\newtheorem{claim}[theorem]{Claim}
\newtheorem{definition}[theorem]{Definition}
\newtheorem{remark}[theorem]{Remark}
\newcommand{\R}{\mathbb R}
\newcommand{\N}{\mathbb N}
\newcommand{\Sn}{\mathbb S}
\newcommand{\Om}{\Omega}
\newcommand{\Sig}{\Sigma}
\newcommand{\eps}{\varepsilon}
\newcommand{\dd}{\,d}
\newcommand{\calH}{\mathcal H}
\newcommand{\calD}{\mathcal D}
\newcommand{\calE}{\mathcal E}
\newcommand{\supp}{\operatorname{supp}}
\newcommand{\diver}{\operatorname{div}}
\newcommand{\Vol}{\operatorname{Vol}}
\newcommand{\dist}{\operatorname{dist}}
\newcommand{\ess}{\operatorname*{ess}}
\newcommand{\HH}{\mathbb H}
\newcommand{\Area}{\operatorname{Area}}
\setcounter{tocdepth}{1}
\allowdisplaybreaks
\begin{document}
\title[Weinstock inequalities for outward-minimizing domains]{Weinstock inequalities for outward-minimizing domains}
\author{Chaoqun Gao, Yong Wei, Rong Zhou}
\date{}
\maketitle
\noindent\emph{Source and licence.} Weinstock inequalities for outward-minimizing domains. Version arXiv:2608.11841v1 (CC BY 4.0). This material is distributed under the Creative Commons Attribution 4.0 International licence, http://creativecommons.org/licenses/by/4.0/, as recorded in the arXiv OAI licence metadata for this exact version and shown on the abstract page https://arxiv.org/abs/2608.11841v1. Full rights notice: licenses/arxiv-2608.11841-CC-BY-4.0.txt.\par\smallskip
\noindent\textit{Editorial scope.} Counted unit: the original \texttt{theorem} environment \texttt{thm:weinstock} at source lines 117--126 together with its complete proof at lines 830--868; the delivered document keeps source lines 83--1040 so that every referenced label and every citation resolves inside it. Native editable source is the paper's own concatenated TeX; the AMS wrapper is editorial formatting only. No publisher class, upstream script or unrelated artwork is redistributed.
\section{Introduction}

Let $(M^n,g)$ be either the Euclidean space $\R^n$ or the hyperbolic space $\HH^n$, and let $\Omega\subset M$ be a bounded connected domain with smooth boundary $\Sigma=\partial\Omega$.  The Steklov eigenvalue problem is
\begin{equation*}
\begin{cases}
    \Delta u=0 & \text{in }\Omega,\\
    \partial_\nu u=\sigma u & \text{on }\Sigma,
\end{cases}
\end{equation*}
where $\Delta$ is the Laplace operator and $\nu$ denotes the outward unit normal.  The first nonzero Steklov eigenvalue is characterized by
\begin{equation}\label{eq:steklov-variational}
    \sigma_1(\Omega)
    =\inf\left\{
        \frac{\int_\Omega |\nabla u|^2\,dv}{\int_{\Sigma} u^2\,d\mu}:
        u\in H^1(\Omega)\setminus\{0\},\quad
        \int_{\Sigma}u\,d\mu=0
      \right\}.
\end{equation}
Here and below $\sigma_1$ denotes the first positive Steklov eigenvalue.  The Steklov problem was introduced by Vladimir Steklov at the turn of the $20$th century.  It is equivalently the spectral problem for the Dirichlet-to-Neumann map, and it arises in inverse problems, hydrodynamics, free boundary minimal surface theory, and differential geometry.  We refer to \cite[Chapter~7]{LMP2023} for an introduction to the Steklov problem and to \cite{CGGS24,Fra20,GP2017} for surveys on recent progress.

One of the central themes in Steklov spectral geometry is the search for isoperimetric upper bounds.  In dimension two, Weinstock \cite{W1954} proved that among simply connected planar domains with prescribed boundary length the disk maximizes $\sigma_1$.  The simple connectivity hypothesis is essential in the plane, and Girouard--Polterovich  \cite[Open Problem~2]{GP2017} ask for the maximal value of $\sigma_1$ among Euclidean domains with prescribed perimeter and for the domains, or limiting sequences of domains, on which it is realized.  In higher Euclidean dimensions, Fraser--Schoen  \cite{FS2019} showed that the direct analogue of Weinstock's theorem is false even among smooth contractible domains.  The recent survey \cite[Section~4]{CGGS24} emphasizes this obstruction and asks for geometric settings in which sharp ball-type inequalities can still be recovered.

Bucur--Ferone--Nitsch--Trombetti \cite{BFNT2021} proved the sharp Euclidean inequality in the class of convex domains.  Kwong--Wei \cite{KwongWei2023} later weakened convexity to the star-shaped mean-convex class by using a weighted three-term isoperimetric inequality \cite{GR2020} with a smooth inverse mean curvature flow argument.  Their theorem fits into the broader program of understanding which geometric assumptions can replace convexity in higher-dimensional Steklov inequalities.  

In space forms, \cite[Open Question~4.27]{CGGS24} asks in particular whether the convex-domain Weinstock inequality has hyperbolic and spherical analogues.  In hyperbolic space, Gu--Li--Wan \cite{GuLiWan2025} proved this for smooth star-shaped mean-convex domains in $\HH^n$, $n\geq4$.  Gu's thesis \cite{GuThesis2026} (see also \cite{GLW26}) contains the additional low-dimensional case $n=3$.  Their proof again uses smooth inverse mean curvature flow, together with a special volume-area ratio function $g$ and an auxiliary function $h$.

The purpose of this paper is to extend both results \cite{KwongWei2023,GuLiWan2025,GuThesis2026} from star-shaped mean-convex domains to outward-minimizing domains.  For a finite-perimeter set $E$ we write $P(E)=\calH^{n-1}(\partial^*E)$.  See \cite[Chapters~12 and~15]{Maggi2012} for the Euclidean theory of finite-perimeter sets and reduced boundaries.  A bounded finite-perimeter set $\Omega\subset M$ is outward minimizing if
\begin{equation*}
    P(E)\geq P(\Omega)
\end{equation*}
for every bounded finite-perimeter set $E$ that contains $\Omega$ up to measure zero.  It is called strictly outward minimizing if equality implies that $E=\Omega$ up to measure zero.  This is a variational perimeter condition rather than a pointwise curvature condition.  In Euclidean space, smooth star-shaped strictly mean-convex domains are strictly outward minimizing \cite[Corollary~3.8]{FM22}, but the outward-minimizing class is much larger and the classical inverse mean curvature flow may develop singularities from such initial data.  We therefore use the weak inverse mean curvature flow of Huisken--Ilmanen \cite{HI2001}.  

Our first theorem is the Euclidean Weinstock inequality in the outward-minimizing class.


\begin{theorem}\label{thm:weinstock}
Let $\Omega\subset\R^n$, $n\geq3$, be a bounded outward-minimizing domain with smooth  boundary $\Sigma=\partial\Omega$.  Then
\begin{equation}\label{eq:weinstock-main}
    \sigma_1(\Omega)|\partial\Omega|^{\frac1{n-1}}
    \leq
    \sigma_1(B^n)|\partial B^n|^{\frac1{n-1}}
    =|\Sn^{n-1}|^{\frac1{n-1}}.
\end{equation}
Equality holds if and only if $\Omega$ is a round ball.
\end{theorem}


A key ingredient for proving Theorem \ref{thm:weinstock} is a three-term geometric inequality in Proposition \ref{prop:weighted-three-term}.  If $r=|x|$ is the Euclidean distance from a fixed origin, it asserts that for every $k>0$
\begin{equation*}
    \int_\Sigma r^k\,d\mu
    \geq
    \frac{n-1}{n-1+k}|\Sn^{n-1}|^{-\frac{k}{n-1}}
    |\Sigma|^{\frac{n-1+k}{n-1}}
    + k\int_\Omega r^{k-1}\,dx.
\end{equation*}
This extends the Euclidean three-term inequality of Kwong--Wei \cite{KwongWei2023} from the smooth star-shaped mean-convex setting to the outward-minimizing setting. Let $\{\Omega_t\}_{t\geq0}$ denote the sublevel-set domains of the weak
inverse mean curvature flow starting from $\Omega$, and set
\[
    \Sigma_t:=\partial^*\Omega_t,
    \qquad
    \mu_t:=\mathcal H^{n-1}\mathbin{\llcorner}\Sigma_t.
\]
The proof is based on the monotonicity of
\begin{equation*}
    |\Sigma_t|^{-\frac{n-1+k}{n-1}}
    \left(\int_{\Sigma_t}r^k\,d\mu_t
        -k\int_{\Omega_t}r^{k-1}\,dx\right)
\end{equation*}
along the weak flow.

Our second theorem is the hyperbolic analogue.

\begin{theorem}\label{thm:hyp-outward-weinstock}
Let $\Omega\subset\HH^n$, $n\geq3$,  be a bounded outward-minimizing domain with smooth boundary $\Sigma=\partial\Omega$.  Let $\Omega^*$ be a geodesic ball satisfying $ |\partial\Omega^*|=|\partial\Omega|$.  Then
\begin{equation}\label{eq:hyp-weinstock-final}
    \sigma_1(\Omega)\leq\sigma_1(\Omega^*).
\end{equation}
Equality holds if and only if $\Omega$ is a geodesic ball.
\end{theorem}

We now describe the proof and the main new points.  In both the Euclidean and hyperbolic arguments, the monotonicity along smooth inverse mean curvature flow is replaced by endpoint distributional monotonicity along the weak flow.  This use of weak IMCF is in the same spirit as its applications to Minkowski-type inequalities for outward-minimizing hypersurfaces.  In Euclidean space, Huisken \cite{Huisken2009Lecture} used weak IMCF to prove the corresponding Minkowski inequality, see also Freire--Schwartz \cite{FS2014}.  Agostiniani--Fogagnolo--Mazzieri \cite{AFM2022} later gave a nonlinear potential theoretic approach.  In hyperbolic space, Harvie \cite{Harvie2026} proved weak-IMCF Minkowski-type inequalities for outward-minimizing domains in dimensions $3\leq n\leq7$.  The measure argument below applies in every dimension $n\geq3$.

The main difficulty is to pass from smooth inverse mean curvature
flow to weak inverse mean curvature flow without losing the
prescribed initial data. After extending $u$ by zero on $\Omega$,
the weak flow admits, for $t>0$, two natural choices of level-set
domains,
\begin{equation*}
    \Omega_t=\{u<t\},
    \qquad
    \Omega_t^+=\operatorname{int}\{u\leq t\}.
\end{equation*}
These sets may differ at a jump time. In particular, the
right-hand limit at $t=0$ may be the strictly minimizing hull
$\Omega_0^+$ rather than the prescribed domain $\Omega$. This
distinction is harmless in some geometric inequalities, but it
matters here because the Steklov estimates are formulated in terms
of the original boundary $\partial\Omega$.

The key device in this paper is the calibrated formulation of weak IMCF \cite[Section~3]{HI2001}.  The calibration is an $L^\infty$ divergence-measure field.  Applying the Gauss--Green formula to a weighted multiple of this field produces a normal trace on the prescribed initial hypersurface.  The trace has the sign needed for monotonicity, and the contribution from the set $\{u=0\}\setminus\overline\Omega$ is included in a nonnegative defect measure.  Thus the distributional identities are written on $[0,\infty)$ with the endpoint value computed on $\Sigma$, not on $\partial\Omega_0^+$. No strict outward-minimizing hypothesis is used.
%This avoids identifying a regularized initial graph across a jump and avoids requiring convergence of the full graph normals in the jump region.

In the Euclidean case, the endpoint transport formula gives the defect-measure identity \eqref{eq:F-defect-measure} behind Proposition \ref{prop:monotonicity}.  In the hyperbolic case, it gives a measure inequality \eqref{eq:hyp-measure-transport} for $G(t)=\int_{\Sigma_t}g\,d\mu_t$.  We combine this inequality with a time-mollification argument for the denominator of the Gu--Li--Wan \cite{GuLiWan2025} functional
\begin{equation}\label{s1.Mt}
    \mathcal M(t)=\frac{1}{|\Sigma_t|}
    \frac{\int_{\Sigma_t}g\,d\mu_t}
        {|\Omega_t|h(\int_{\Omega_t}\frac{\lambda'g}{\lambda}\,dv)}.
\end{equation}
The weighted ball integral is convex as a function of the ball volume. Jensen's inequality therefore preserves the mass-transplantation estimate under time mollification, while the
measure formulation automatically includes all jump contributions. This yields monotonicity of $ \mathcal M(t)$ along weak IMCF in hyperbolic space. The remaining hyperbolic estimates are the radial mass-transplantation inequalities of Gu--Li--Wan \cite{GuLiWan2025} and the low-dimensional comparison from Gu's thesis \cite{GuThesis2026}.


The paper is organized as follows.  Section \ref{sec.2} recalls the weak inverse mean curvature flow facts used later and develops the calibrated endpoint transport formula.  Section \ref{sec.3} proves the Euclidean weighted three-term inequality and then Theorem \ref{thm:weinstock}.  Section \ref{sec.4} proves the hyperbolic weak-flow monotonicity and then Theorem \ref{thm:hyp-outward-weinstock}.

\section{Weak inverse mean curvature flow}\label{sec.2}

We recall the weak inverse mean curvature flow in the form needed below.  Huisken--Ilmanen \cite{HI2001} formulated the weak flow on general ambient Riemannian manifolds.  In this paper the ambient manifold $M$ is either $\R^n$ or $\HH^n$.  For a finite-perimeter set $E$ we write $P(E)=\calH^{n-1}(\partial^*E)$. This is the reduced-boundary representation of the perimeter.  See \cite[Theorem~15.9]{Maggi2012}.  All integrals over weak level-set boundaries are taken over the reduced boundary.

Let $\Omega\subset M$ be a bounded domain with smooth boundary $\Sigma$.  A proper function
\begin{equation*}
    u\in C^0(M\setminus\Omega)
      \cap W_{\mathrm{loc}}^{1,\infty}(M\setminus\overline\Omega)
\end{equation*}
is a weak solution of inverse mean curvature flow with initial data $\Sigma$ if its continuous boundary trace is zero on $\Sigma$, if $u\to\infty$ at infinity, and if the sublevel sets
\begin{equation*}
    \Omega_t:=\Omega\cup\{x\in M\setminus\overline\Omega:u(x)<t\},
    \qquad
    \Sigma_t:=\partial^*\Omega_t
\end{equation*}
solve
\begin{equation}\label{eq:level-set-imcf}
    \diver_M\left(\frac{\nabla u}{|\nabla u|}\right)=|\nabla u|
\end{equation}
in the variational sense of Huisken--Ilmanen.  If $u$ is smooth and $|\nabla u|\neq0$, then \eqref{eq:level-set-imcf} is equivalent to the classical inverse mean curvature flow \cite{Gerhardt1990,Gerhardt2011,Urbas1990}
\begin{equation*}
    \partial_t X=\frac{1}{H}\nu.
\end{equation*}

The explicit expanding-sphere solutions in $\R^n$ and $\HH^n$ are proper weak subsolutions.  The weak existence theorem \cite[Theorem~3.1]{HI2001} therefore gives a proper weak solution for every smooth bounded initial domain in the settings used here. The solution obtained by this construction satisfies
$u\geq0$ on $M\setminus\Omega$. Moreover, the boundary gradient estimate in the proof of \cite[Theorem~3.1]{HI2001}, together with the continuous zero trace of $u$ on $\Sigma$, implies that the extension obtained by setting $u=0$ on $\Omega$ belongs to $W_{\mathrm{loc}}^{1,\infty}(M)$. We continue to denote this extension by $u$. For $t>0$, the sets $\Omega_t$ are outward minimizing.  The right-continuous representatives
\begin{equation*}
    \Omega_t^+:=\operatorname{int}\bigl(\Omega\cup\{x\in M\setminus\overline\Omega:u(x)\leq t\}\bigr)
\end{equation*}
are strictly outward minimizing by the minimizing hull property \cite[Minimizing Hull Property~1.4]{HI2001}.  A jump occurs when $\Omega_t$ and $\Omega_t^+$ differ by a set of positive measure.

The monotonicity of the sublevel sets gives
\begin{equation*}
    \chi_{\Omega_s}\longrightarrow\chi_{\Omega_t}
    \quad\text{in }L^1_{\mathrm{loc}}\text{ as }s\nearrow t,
\end{equation*}
and
\begin{equation*}
    \chi_{\Omega_s}\longrightarrow\chi_{\Omega_t^+}
    \quad\text{in }L^1_{\mathrm{loc}}\text{ as }s\searrow t.
\end{equation*}
At $t=0$, the right limit may be the strictly minimizing hull $\Omega_0^+$ rather than the prescribed domain $\Omega$.  We set $\Omega_0=\Omega$ and $\Sigma_0=\Sigma$ throughout.  Every distributional identity below contains an explicit endpoint term at $t=0$ computed on the prescribed initial domain.

If the initial domain is outward minimizing, the exponential area law \cite[Lemma~1.6]{HI2001} gives
\begin{equation}\label{eq:area-growth}
    P(\Omega_t)=e^tP(\Omega),
    \qquad t\geq0.
\end{equation}
The equality at $t=0$ uses the prescribed initial trace.  We write $|\Sigma_t|=P(\Omega_t)$.

The proper weak solution obtained through the elliptic-regularization construction in \cite[Theorem 3.1]{HI2001} carries the corresponding calibration. More precisely, there exists a measurable vector field $\nu$ on $M\setminus\overline\Omega$ such that (see \cite[(3.15)]{HI2001})
\begin{equation}\label{eq:calibration}
    |\nu|\leq1,
    \qquad
    \langle\nabla u,\nu\rangle=|\nabla u|
    \quad\text{a.e. in }M\setminus\overline\Omega,
\end{equation}
and
\begin{equation}\label{eq:calibration-div}
    \int_{M\setminus\overline\Omega}\langle\nabla\phi,\nu\rangle\,dv
    =-\int_{M\setminus\overline\Omega}\phi|\nabla u|\,dv,
    \qquad
    \phi\in C_c^1(M\setminus\overline\Omega).
\end{equation}
Thus $\diver\nu=|\nabla u|$ in distributions.  At regular points of a level set, $\nu=\nabla u/|\nabla u|$ is the outward unit normal.

We use one endpoint convention throughout the paper.  If $\Phi\in L^1_{\mathrm{loc}}([0,\infty))$ has the prescribed value $\Phi(0)$, let $D_t\Phi$ denote the distributional derivative on $\mathbb R$ of its constant extension to $(-\infty,0)$, obtained by setting $\Phi(t)=\Phi(0)$ for $t<0$.  Then every $\eta\in C_c^1(\mathbb R)$ satisfies
\begin{equation}\label{eq:endpoint-derivative-convention}
    \langle D_t\Phi,\eta\rangle
    =-\int_0^\infty\Phi(t)\eta'(t)\,dt-\eta(0)\Phi(0).
\end{equation}
Here and below, $\eta\in C_c^1([0,\infty))$ means that $\eta$ is the restriction to $[0,\infty)$ of some function $\widetilde{\eta}\in C_c^1(\mathbb R)$. In particular, $\eta(0)$ is not required to vanish. If, in addition, the right trace $\Phi(0+)$ exists and $D_t\Phi$ is a Radon measure, then
\begin{equation*}
    D_t\Phi(\{0\})=\Phi(0+)-\Phi(0).
\end{equation*}
Thus the difference between the prescribed endpoint value and the positive-time right trace is explicitly retained at $t=0$.

All distributional identities obtained from this convention are identities on $\mathbb R$.  When a density such as $A(t)\,dt$ is defined only for $t>0$, it means $A(t)\mathbf 1_{(0,\infty)}(t)\,dt$ on $\mathbb R$.

\subsection{An endpoint Gauss--Green formula}\label{subsec:endpoint-GG}

The next lemma is the basic endpoint identity used in both ambient geometries.  The normal trace is understood in the sense of divergence-measure fields.  We use the Gauss--Green formula and the normal trace from \cite[Theorem~2.2]{ChenFrid1999}, and the product rule from \cite[Theorem~3.1]{ChenFrid1999}.  These results are local, and the coordinate reduction below gives the corresponding formula on a smooth
Riemannian manifold.  For a recent intrinsic formulation on regular
domains in metric measure spaces, including the Gauss--Green formula
and the sharp normal-trace estimate, see \cite[Theorem~3.6]{GornyMazon2024}.

\begin{lemma}\label{lem:endpoint-transport}
Let $f\in W^{1,\infty}_{\mathrm{loc}}(M)$ be nonnegative.  For almost every $t>0$ set
\begin{equation*}
    A_f(t):=\int_{\partial^*\Omega_t}f\,d\calH^{n-1},
    \qquad
    A_f(0):=\int_\Sigma f\,d\mu.
\end{equation*}
Then $A_f\in L^1_{\mathrm{loc}}([0,\infty))$.  There is a normal trace
\begin{equation*}
    \tau:=\operatorname{Tr}_\Sigma(\nu\cdot\nu_\Omega),
    \qquad |\tau|\leq1,
\end{equation*}
such that every $\eta\in C_c^1([0,\infty))$ satisfies
\begin{align}
    -\int_0^\infty\eta'(t)A_f(t)\,dt-\eta(0)A_f(0)
    =&\int_0^\infty\eta(t)A_f(t)\,dt 
      +\int_{M\setminus\overline\Omega}\eta(u)\langle\nabla f,\nu\rangle\,dv \notag\\
    &+\eta(0)\int_\Sigma f(\tau-1)\,d\mu.\label{eq:endpoint-transport-exact}
\end{align}
Consequently, if $\eta\geq0$, then
\begin{equation}\label{eq:endpoint-transport-ineq}
    -\int_0^\infty\eta'(t)A_f(t)\,dt-\eta(0)A_f(0)
    \leq
    \int_0^\infty\eta(t)A_f(t)\,dt
    +\int_{M\setminus\overline\Omega}\eta(u)\langle\nabla f,\nu\rangle\,dv.
\end{equation}
\end{lemma}

\begin{proof}
Since $u$ is proper and $\eta$ has compact support, the vector field
\begin{equation*}
    X:=f\eta(u)\nu
\end{equation*}
has compact support in the closure of $M\setminus\overline\Omega$.  Equations \eqref{eq:calibration} and \eqref{eq:calibration-div}, together with the product rule for divergence-measure fields, give
\begin{equation}
    \diver X
    =\eta(u)\langle\nabla f,\nu\rangle
     +f\bigl(\eta'(u)+\eta(u)\bigr)|\nabla u|
     \label{equ-divX}
\end{equation}
as an identity of distributions on $M\setminus\overline\Omega$.  In particular, $\diver X$ is represented there by a locally integrable function.

We next explain why the Gauss--Green formula applies in the Riemannian
setting. Choose a bounded smooth open set $U$ containing $\overline\Omega$ and the support of $X$, with $X=0$ near $\partial U$.  In a coordinate chart $y=(y^1,\ldots,y^n)$, write $X=X^i\partial_i$ and define the Euclidean coordinate field
\begin{equation*}
    Y^i:=\sqrt{\det(g_{ab})}\,X^i.
\end{equation*}
Then
\begin{equation*}
    \diver_{\mathbb R^n}Y\,dy
    =\diver_gX\,dv_g.
\end{equation*}
Thus the Euclidean divergence-measure theory in
\cite[Theorems~2.2 and~3.1]{ChenFrid1999} applies to $Y$ in every chart.  Under this coordinate change, its boundary flux becomes the intrinsic Riemannian flux $\langle X,\nu_E\rangle_g\,d\calH_g^{n-1}$, where $\nu_E$ is the measure-theoretic outward unit normal of the set under consideration.  A finite coordinate cover and a partition of unity therefore give the Gauss--Green formula on $U\setminus\overline\Omega$.

The same coordinate argument applied to the calibration field $\nu$ gives its normal trace on $\Sigma$.  This does not assert that $\nu$ has a pointwise boundary value. It means that the flux of $\nu$ through $\Sigma$, taken from the side $M\setminus\overline\Omega$, is well defined by the Gauss--Green formula.  With $\nu_{\Omega}$ denoting the unit normal pointing out of $\Omega$, write
\begin{equation*}
    \tau=\operatorname{Tr}_\Sigma(\nu\cdot\nu_\Omega).
\end{equation*}

Using the normal exponential deformation of $\Sigma$ in the weak-star characterization of the normal
trace in \cite[Theorem~2.2(iii)]{ChenFrid1999} and the bound
$|\nu|\leq1$, we obtain
\begin{equation}\label{eq:normal-trace-bound}
    \|\tau\|_{L^\infty(\Sigma)}
    \leq
    \|\nu\|_{L^\infty(U\setminus\overline\Omega)}
    \leq1.
\end{equation}
Since $\eta(u)$ is locally Lipschitz and has trace
$\eta(0)$ on $\Sigma$, the product rule and the characterization of the normal trace therefore give
\begin{equation*}
    \operatorname{Tr}_\Sigma(X\cdot\nu_\Omega)
    =
    f\eta(0)\operatorname{Tr}_\Sigma(\nu\cdot\nu_\Omega)
    =
    f\eta(0)\tau.
\end{equation*} 
The flux across $\partial U$ is zero.  Since the outer unit normal of $U\setminus\overline\Omega$ along $\Sigma$ is $-\nu_\Omega$, the Gauss--Green formula yields
\begin{equation}\label{s2.pf1}
    \int_{M\setminus\overline\Omega}\diver X\,dv
    =-\eta(0)\int_\Sigma f\tau\,d\mu.
\end{equation}

For each $T>0$, properness of $u$ implies that $\{0\leq u\leq T\}$ is compact. Since $f$ and $|\nabla u|$ are locally bounded, the coarea formula gives
\begin{equation*}
    \int_0^T A_f(t) \dd t = \int_{\{0<u<T\}} f |\nabla u| \dd v<\infty.
\end{equation*}
Consequently $A_f\in L_{\mathrm{loc}}^1([0,\infty))$. Since $u\in W_{\mathrm{loc}}^{1,\infty}(M)$ and
$\Omega_t=\{u<t\}$ for every $t>0$, applying the coarea formula to \eqref{equ-divX} gives
\begin{align}\label{s2.pf2}
 \int_{M\setminus\overline\Omega}\diver X\,dv
   =&  \int_{M\setminus\overline\Omega}\eta(u)\langle\nabla f,\nu\rangle \,dv+
    \int_{M\setminus\overline\Omega}
        f\bigl(\eta'(u)+\eta(u)\bigr)|\nabla u|\,dv\nonumber\\
    =&\int_{M\setminus\overline\Omega}\eta(u)\langle\nabla f,\nu\rangle \,dv+
    \int_0^\infty
        (\eta'(t)+\eta(t))
        \left(
            \int_{\partial^*\Omega_t}
                f\,d\mathcal H^{n-1}
        \right)dt \nonumber\\
    =&\int_{M\setminus\overline\Omega}\eta(u)\langle\nabla f,\nu\rangle \,dv+
    \int_0^\infty
        (\eta'(t)+\eta(t))A_f(t)\,dt.
\end{align}
Properness of $u$ and the local boundedness of $f$ ensure that all the integrals are finite. The set $\{u=0\}\setminus\overline\Omega$ causes no additional coarea term, because $|\nabla u|=0$ almost everywhere on this set. The coarea formula determines $A_f(t)$ only for almost every $t>0$. It does not determine the separately prescribed endpoint value $A_f(0)=\int_\Sigma f\,d\mu$.

Combining \eqref{s2.pf1} and \eqref{s2.pf2} and rearranging gives \eqref{eq:endpoint-transport-exact}.  Since $f\geq0$ and $\tau\leq1$ by \eqref{eq:normal-trace-bound}, the boundary defect term is non-positive
for every nonnegative $\eta$.  Hence
\eqref{eq:endpoint-transport-ineq} follows.
\end{proof}


\section{The Euclidean case}\label{sec.3}

In this section, we prove the Euclidean Weinstock inequality.  The first step is a weak-flow proof of the weighted three-term inequality.  We apply the endpoint Gauss--Green formula from Lemma \ref{lem:endpoint-transport}, combine it with a finite-perimeter divergence estimate, and obtain monotonicity of a normalized quantity. The final subsection applies the case $k=1$ to the Steklov variational characterization.

Throughout this section, we fix an origin and let $r(x):=|x|$ denote the Euclidean distance from it. Away from the origin, we write $\partial_r:=\nabla r=\frac{x}{r}$. Whenever $\partial_r$ occurs in an almost-everywhere identity, its value at the origin may be chosen arbitrarily. We reserve $D$ for
distributional derivatives.

\subsection{The weighted three-term inequality}

\begin{proposition}\label{prop:weighted-three-term}
Let $\Om\subset\R^n$, $n\geq3$, be a bounded outward-minimizing domain with smooth boundary $\Sig=\partial\Om$.  Fix an origin in $\R^n$ and set $r=|x|$.  Then for every $k>0$,
\begin{equation}\label{eq:weighted-main}
    \int_\Sig r^k\,d\mu
    \geq
    \frac{n-1}{n-1+k}|\Sn^{n-1}|^{-\frac{k}{n-1}}
    |\Sig|^{\frac{n-1+k}{n-1}}
    + k\int_\Om r^{k-1}\,dx.
\end{equation}
Equality holds if and only if $\Sig$ is a round sphere centered at the chosen origin.
\end{proposition}

Fix $k>0$. For every $t>0$, set
\begin{equation*}
    B_k(t):=k\int_{\Om_t}r^{k-1}\,dx,
\end{equation*}
and, for almost every $t>0$, set
\begin{equation*}
    A_k(t):=\int_{\Sig_t}r^k\,d\mu_t,
    \qquad
    F_k(t):=A_k(t)-B_k(t).
\end{equation*}
At $t=0$, prescribe
\begin{equation*}
    A_k(0):=\int_{\Sig}r^k\,d\mu,
    \qquad
    B_k(0):=k\int_{\Om}r^{k-1}\,dx,
    \qquad
    F_k(0):=A_k(0)-B_k(0).
\end{equation*}
Whenever distributional derivatives in time are used, these functions are extended to $t<0$ by their prescribed values at $t=0$.

\begin{lemma}\label{lem:transport}
Let $\{\Om_t\}$ be the weak inverse mean curvature flow starting from a smooth bounded domain $\Om$.  Then $F_k\in L^1_{\mathrm{loc}}([0,\infty))$ and
\begin{equation*}
    D_tF_k\leq A_k\,dt
\end{equation*}
as an inequality of distributions on $\mathbb R$ under the endpoint convention above. Here $A_k\,dt$ is zero on $(-\infty,0)$.  Equivalently, for every nonnegative $\eta\in C_c^1([0,\infty))$,
\begin{equation}\label{eq:distribution-F}
    -\int_0^\infty\eta'(t)F_k(t)\,dt-\eta(0)F_k(0)
    \leq
    \int_0^\infty\eta(t)A_k(t)\,dt.
\end{equation}
More precisely, define
\begin{equation}\label{eq:Euclidean-defect-measure-definition}
\begin{split}
    \mathfrak e_k
    :=&{u}_{\#}\left(
      kr^{k-1}\bigl(1-\partial_r\cdot\nu\bigr)
      \mathcal L^n\mathbin{\llcorner}
      (\R^n\setminus\overline\Om)\right)\\
    &+
      \left(\int_\Sig r^k(1-\tau)\,d\mu\right)\delta_0,
\end{split}
\end{equation}
where $\tau$ is the normal trace from Lemma \ref{lem:endpoint-transport}.  This is a nonnegative Radon measure and
\begin{equation}\label{eq:F-defect-measure}
    D_tF_k=A_k\,dt-\mathfrak e_k.
\end{equation}
Equivalently, for every $\eta\in C_c([0,\infty))$,
\begin{align}
    \int_{[0,\infty)}\eta\,d\mathfrak e_k
    &=\int_{\R^n\setminus\overline\Om}
      \eta(u)kr^{k-1}\bigl(1-\partial_r\cdot\nu\bigr)\,dx \notag\\
    &\quad+\eta(0)\int_\Sig r^k(1-\tau)\,d\mu.\label{eq:Euclidean-defect}
\end{align}
\end{lemma}

\begin{proof}
For $\delta>0$ define
\begin{equation*}
    f_\delta(x):=(r^2+\delta^2)^{k/2}.
\end{equation*}
The regularization removes the possible singularity of $\nabla r^k$ at the origin and allows a single argument for every $k>0$.  The function $f_\delta$ is smooth and nonnegative, so Lemma \ref{lem:endpoint-transport} applies to it.  We now justify the passage to the limit $\delta\downarrow0$.  Fix $\eta$ and choose $T>0$ such that the supports of $\eta$ and $\eta'$ are contained in $[0,T]$.  By the properness of $u$, the set
\begin{equation*}
    K_T:=\overline{\{x\in\R^n\setminus\overline\Om:0\leq u(x)\leq T\}}
\end{equation*}
is compact.  For $\zeta=\eta$ and $\zeta=\eta'$, the coarea formula gives
\begin{equation*}
    \int_0^\infty \zeta(t)A_{f_\delta}(t)\,dt
    =\int_{\R^n\setminus\overline\Om}
      \zeta(u)f_\delta|\nabla u|\,dx.
\end{equation*}
Since $f_\delta\to r^k$ uniformly on $K_T$ and $|\nabla u|$ is locally integrable, these two terms converge to the corresponding terms with $r^k$.  The endpoint term also converges because $f_\delta\to r^k$ uniformly on the smooth compact hypersurface $\Sig$.

It remains to pass to the limit in the gradient term.  On $K_T$,
\begin{equation*}
    \nabla f_\delta
    =kr(r^2+\delta^2)^{\frac{k}{2}-1}\partial_r
    \longrightarrow kr^{k-1}\partial_r
\end{equation*}
almost everywhere.  If $0<k<2$, then $|\nabla f_\delta|\leq kr^{k-1}$ away from $\{0\}$, and $r^{k-1}$ is locally integrable for every $k>0$.  If $k\geq2$, then $|\nabla f_\delta|$ is uniformly bounded on $K_T$.  Since $|\nu|\leq1$, dominated convergence gives
\begin{equation*}
    \int_{\R^n\setminus\overline\Om}
      \eta(u)\langle\nabla f_\delta,\nu\rangle\,dx
    \longrightarrow
    \int_{\R^n\setminus\overline\Om}
      \eta(u)kr^{k-1}\partial_r\cdot\nu\,dx.
\end{equation*}
Thus we may let $\delta\downarrow0$ in \eqref{eq:endpoint-transport-exact}. This gives
\begin{align}
    -\int_0^\infty\eta'(t)A_k(t)\,dt-\eta(0)A_k(0)
    &=\int_0^\infty\eta(t)A_k(t)\,dt 
      +\int_{\R^n\setminus\overline\Om}
       \eta(u)kr^{k-1}\partial_r\cdot\nu\,dx \notag\\
    &\quad+\eta(0)\int_\Sig r^k(\tau-1)\,d\mu.\label{eq:Ak-endpoint}
\end{align}

Since $u$ is proper and $kr(x)^{k-1}$ is locally integrable,
\begin{equation*}
    \int_0^\infty\int_{\mathbb R^n\setminus\Omega}
    |\eta'(t)|\mathbf 1_{\{u(x)<t\}}kr(x)^{k-1}\,dx\,dt
    \leq
    T\|\eta'\|_{L^\infty}\int_{K_T}kr(x)^{k-1}\,dx<\infty.
\end{equation*}
Hence Fubini's theorem applies. Since $\eta$ has compact support and
\begin{equation*}
    B_k(t)=B_k(0)
      +\int_{\R^n\setminus\overline\Om}
        \mathbf 1_{\{u<t\}}kr^{k-1}\,dx,
\end{equation*}
we have 
\begin{align}\label{eq:Bk-distribution}
 -\int_0^\infty\eta'(t)&B_k(t)\,dt-\eta(0)B_k(0)\nonumber\\
 =&-\int_0^\infty\eta'(t)\left(B_k(0)
      +\int_{\R^n\setminus\overline\Om}
        \mathbf 1_{\{u<t\}}kr^{k-1}\,dx\right)\,dt-\eta(0)B_k(0)\nonumber\\
=&-B_k(0)\int_0^\infty\eta'(t)dt-\eta(0)B_k(0)\nonumber\\
&\qquad -\int_0^\infty\eta'(t)\int_{\R^n\setminus\overline\Om}
        \mathbf 1_{\{u<t\}}kr^{k-1}\,dx \,dt\nonumber\\
 =&-\int_{\mathbb R^n\setminus\Omega}kr^{k-1}
   \int_0^\infty\eta'(t)\mathbf 1_{\{u(x)<t\}}\,dt\,dx\nonumber\\
 =&-\int_{\mathbb R^n\setminus\Omega}kr^{k-1}
   \int_{u(x)}^\infty\eta'(t)\,dt\,dx\nonumber\\
 =&\int_{\mathbb R^n\setminus\Omega}
   \eta(u)kr^{k-1}\,dx.
\end{align}
Subtracting \eqref{eq:Bk-distribution} from \eqref{eq:Ak-endpoint} gives \eqref{eq:F-defect-measure} with the measure in \eqref{eq:Euclidean-defect-measure-definition}.   The identity first holds for test functions in $C_c^1([0,\infty))$ and then for every function in $C_c([0,\infty))$ by density.  The measure is nonnegative because $|\nu|\leq1$ and $\tau\leq1$.  Equation \eqref{eq:distribution-F} follows.
\end{proof}

We next relate $A_k(t)$ and $B_k(t)$ by a divergence estimate.

\begin{lemma}\label{lem:divergence}
Let $E\subset\R^n$ be a finite-perimeter set with compact closure.  For every $k>0$,
\begin{equation}\label{eq:divergence-estimate}
    (n-1+k)\int_E r^{k-1}\,dx
    \leq
    \int_{\partial^*E}r^k\,d\calH^{n-1}.
\end{equation}
If $|E|>0$, equality holds if and only if $E$ agrees almost everywhere with a ball centered at the chosen origin. In particular, for almost every $t>0$ along the weak flow,
\begin{equation}\label{eq:B-vs-A}
   % B_k(t)\leq\frac{k}{n-1+k}A_k(t),
   % \qquad
    F_k(t)\geq\frac{n-1}{n-1+k}A_k(t).
\end{equation}
\end{lemma}

\begin{proof}
Consider the smooth vector fields
\begin{equation*}
    X_\delta=(r^2+\delta^2)^{\frac{k-1}{2}}x.
\end{equation*}
The limiting field $r^{k-1}x$ has the desired divergence but may fail to be $C^1$ at the origin when $0<k<1$.  The field $X_\delta$ regularizes this singularity.  Since $E$ has compact closure, we may multiply $X_\delta$ by a smooth cutoff that equals one near $\overline E$.  The Gauss--Green theorem for finite-perimeter sets \cite[Theorem~15.9]{Maggi2012} then yields
\begin{equation}\label{s3.delta}
\int_E\diver X_\delta\,dx
=\int_{\partial^*E} X_\delta\cdot\nu_E\,d\calH^{n-1}.
\end{equation}
A direct computation gives
\begin{equation*}
    \diver X_{\delta} = n \left( r^2 + \delta^2 \right)^{\frac{k-1}{2}} + (k-1)\left( r^2 + \delta^2 \right)^{\frac{k-3}{2}}r^2.
\end{equation*}
Passing to the limit as $\delta\downarrow 0$ in \eqref{s3.delta} gives
\begin{equation}\label{eq:weighted-divergence-identity}
    (n-1+k)\int_E r^{k-1}\,dx
    =\int_{\partial^*E}r^{k-1}x\cdot\nu_E\,d\calH^{n-1}.
\end{equation}
Indeed, if $0<k<1$, then
\begin{equation*}
    |\diver X_\delta|  \leq (n+1) r^{k-1},
    \qquad
    |X_\delta|\leq r^k,
\end{equation*}
and $r^{k-1}$ is locally integrable in $\mathbb{R}^n$, while for $k\geq1$ both quantities are locally bounded uniformly in $\delta$.  Dominated convergence therefore applies to the volume term and the reduced-boundary term in \eqref{s3.delta}. Since $x\cdot\nu_E\leq r$, inequality \eqref{eq:divergence-estimate} follows.


Suppose equality holds and $|E|>0$.  Equation \eqref{eq:weighted-divergence-identity} implies
\begin{equation*}
    \nu_E=\partial_r
    \quad\text{for }\calH^{n-1}\text{-almost every point of }\partial^*E.
\end{equation*}
Indeed, the nonnegative integrand $r^{k-1}(r-x\cdot\nu_E)$ has zero integral. Recall that
\begin{equation*}
    D\chi_E=-\nu_E|D\chi_E|.
\end{equation*}
Then we have $D\chi_E=-\partial_r|D\chi_E|$, where $|D\chi_E|=\mathcal{H}^{n-1}\mathbin{\llcorner}\partial^*E$.  The value of $\partial_r$ at the origin is immaterial. We first prove that $E$ is rotationally invariant. Let $A$ be a skew-symmetric matrix and set $Z_A(x)=Ax$.  Since 
\begin{equation*}
    Z_A\cdot \partial_r=\frac{1}{r}Ax\cdot x=\frac{1}{r}x^TA^Tx=0
\end{equation*}
and $\diver Z_A=\operatorname{tr}A=0$,
\begin{equation*}
    \diver(\chi_E Z_A)
    =Z_A\cdot D\chi_E+\chi_E\diver Z_A=0
\end{equation*}
in the sense of distributions. If $R_t=e^{tA}$, then for every $\varphi\in C_c^1(\R^n)$,
\begin{align*}
\frac{d}{dt}\int_{\mathbb R^n}
 \chi_E(x)\varphi(R_{-t}x)\,dx
&=-\int_{\mathbb R^n}
 \chi_E(x)Z_A(x)\cdot
 \nabla(\varphi\circ R_{-t})(x)\,dx\\
&=\left\langle
 \diver(\chi_EZ_A),\varphi\circ R_{-t}
 \right\rangle=0.
\end{align*}
Thus the integral is independent of $t$. By a change of variables $x=R_ty$ and noting that $\det R_t=1$, 
\begin{equation*}
    \int_{\R^n}\chi_E(R_ty)\varphi(y)dy=\int_{\R^n}\chi_E(y)\varphi(y)dy
\end{equation*}
for every $\varphi\in C_c^1(\R^n)$. This yields $\chi_E\circ R_t=\chi_E$ almost everywhere. Since $A$ was arbitrary and the rotations $e^{tA}$ generate $\mathrm{SO}(n)$, the function $\chi_E$ is invariant almost everywhere under every rotation about the origin. A standard averaging argument
over $\mathrm{SO}(n)$, together with Fubini's theorem, then shows that $\chi_E$ agrees almost everywhere with a radial function. Since $\chi_E$ takes only the values zero and one, there exists a measurable function $\gamma:(0,\infty)\to\{0,1\}$ such that $\chi_E(x)=\gamma(|x|)$ almost everywhere.

To determine the radial profile, let $0\leq\psi\in C_c^1((0,\infty))$ and define
\begin{equation*}
    Y(x)=\frac{\psi(r)}{r^{n-1}}\partial_r.
\end{equation*}
Then $Y$ is smooth and compactly supported away from the origin, and $\diver Y=r^{1-n}\psi'(r)$.  The Gauss--Green formula and $\nu_E=\partial_r$ give
\begin{align*}
    |\Sn^{n-1}|\int_0^\infty \gamma(r)\psi'(r)\,dr
    &=\int_E\diver Y\,dx\\
    &=\int_{\partial^*E}\frac{\psi(r)}{r^{n-1}}\,d\calH^{n-1}\geq0.
\end{align*}
Thus $D_r\gamma\leq0$ in distributions, so $\gamma$ has a non-increasing representative.  Since $\gamma$ takes only the values zero and one, compactness and positive measure imply
\begin{equation*}
    \chi_E=\mathbf 1_{B_R}
\end{equation*}
almost everywhere for some $R>0$.  The converse follows by direct computation.  Applying \eqref{eq:divergence-estimate} to $E=\Om_t$ gives \eqref{eq:B-vs-A} for almost every $t>0$.
\end{proof}

\begin{proposition}[Monotonicity]\label{prop:monotonicity}
Let $\Om\subset\R^n$ be a smooth outward-minimizing domain, and let $\{\Om_t\}$ be the weak inverse mean curvature flow starting from $\Om$. For every $k>0$, define, for almost every $t>0$ and at the prescribed value $t=0$,
\begin{equation*}
    Q_k(t):=|\Sig_t|^{-\frac{n-1+k}{n-1}}F_k(t).
\end{equation*}
Then $Q_k$ agrees almost everywhere on $(0,\infty)$ with a non-increasing function, whose value at $t=0$ is the prescribed value computed on the initial boundary.
\end{proposition}

\begin{proof}
First define $H_k(t):=e^{-\frac{n-1+k}{n-1} t}F_k(t)$ for $t\geq0$, with $H_k(0)=F_k(0)$, and then extend $H_k$ itself constantly to $t<0$. Equivalently, $H_k=wF_k$ on $\mathbb R$, where $w(t)=1$ for $t<0$ and $w(t)=e^{-\frac{n-1+k}{n-1} t}$ for $t\geq0$. The product rule for BV functions, \eqref{eq:F-defect-measure}, and \eqref{eq:B-vs-A} give
\begin{equation}\label{eq:Q-defect}
    D_tH_k
    =-e^{-\frac{n-1+k}{n-1} t}\left(
      \mathfrak e_k+\left(\frac{n-1+k}{n-1} F_k-A_k\right)\,dt\right).
\end{equation}
The right-hand side is a nonpositive Radon measure. Hence $H_k$ agrees almost everywhere with a non-increasing function. Since $|\Sig_t|=e^t|\Sig|$ by \eqref{eq:area-growth},
\begin{equation*}
    Q_k(t)=|\Sig|^{-\frac{n-1+k}{n-1}}H_k(t)
\end{equation*}
for almost every $t>0$ and at $t=0$.  This gives the asserted representative of $Q_k$.
\end{proof}

\begin{proof}[Proof of Proposition \ref{prop:weighted-three-term}]

We first compute the limit of $Q_k(t)$ as $t\to\infty$. By the eventual smoothness theorem of Huisken--Ilmanen \cite[Theorem~2.7]{HI2008}, after increasing $T>0$ if necessary, the weak flow agrees for $t\geq T$ with a smooth star-shaped mean-convex inverse mean curvature flow. Set
\begin{equation*}
    \widetilde\Omega_t:=e^{-\frac{t}{n-1}}\Omega_t,
    \qquad
    \widetilde\Sigma_t:=e^{-\frac{t}{n-1}}\Sigma_t .
\end{equation*}
The asymptotic convergence of the rescaled inverse mean curvature flow \cite{Gerhardt1990,Urbas1990} implies that, for all sufficiently large $t$, one may write
\begin{equation*}
    \widetilde\Sigma_t
    =\{\rho_t(\theta)\theta:\theta\in\mathbb S^{n-1}\},
    \qquad
    \rho_t\longrightarrow R_\infty
    \quad\text{in }C^1(\mathbb S^{n-1}),
\end{equation*}
where, by the area law \eqref{eq:area-growth},
\begin{equation*}
    R_\infty
    =\left(\frac{|\Sigma|}{|\mathbb S^{n-1}|}\right)^{\frac1{n-1}}.
\end{equation*}
For every $k>0$, polar coordinates give
\begin{equation*}
    k\int_{\widetilde\Omega_t}r^{k-1}\,dx
    =\frac{k}{n-1+k}
      \int_{\mathbb S^{n-1}}\rho_t^{\,n-1+k}\,d\theta,
\end{equation*}
while the radial graph representation gives
\begin{equation*}
    \int_{\widetilde\Sigma_t}r^k\,d\widetilde\mu_t
    =
    \int_{\mathbb S^{n-1}}
       \rho_t^{\,n-1+k}
       \sqrt{1+\frac{|\overline{\nabla}\rho_t|_{g_{\mathbb{S}^{n-1}}}^2}{\rho_t^2}}\,d\theta,
\end{equation*}
where $\overline{\nabla}$ denotes the Levi-Civita connection on $\mathbb{S}^{n-1}$. Consequently,
\begin{equation*}
    \lim_{t\to\infty}
    \left(
      \int_{\widetilde\Sigma_t}r^k\,d\widetilde\mu_t
      -k\int_{\widetilde\Omega_t}r^{k-1}\,dx
    \right)
    =
    \frac{n-1}{n-1+k}
    |\mathbb S^{n-1}|R_\infty^{\,n-1+k}.
\end{equation*}
Since $Q_k$ is invariant under Euclidean dilations and
\begin{equation*}
    |\widetilde\Sigma_t|
    =|\Sigma|
    =|\mathbb S^{n-1}|R_\infty^{\,n-1},
\end{equation*}
we obtain
\begin{equation}\label{eq:limit-Q}
    \lim_{t\to\infty}Q_k(t)
    =
    \frac{n-1}{n-1+k}
    |\mathbb S^{n-1}|^{-\frac{k}{n-1}}.
\end{equation}
This argument applies to every $k>0$. In particular, the possible singularity of $r^{k-1}$ at the origin when $0<k<1$ is harmless in the polar-coordinate formula above.

By Proposition \ref{prop:monotonicity} and \eqref{eq:limit-Q},
\begin{equation*}
    |\Sig|^{-\frac{n-1+k}{n-1}}
    \left(\int_\Sig r^k\,d\mu-k\int_\Om r^{k-1}\,dx\right)
    =Q_k(0)
    \geq
    \frac{n-1}{n-1+k}|\Sn^{n-1}|^{-\frac{k}{n-1}}.
\end{equation*}
Rearranging gives \eqref{eq:weighted-main}.

Suppose that equality holds, and let $Q_k$ be the non-increasing
representative provided by Proposition \ref{prop:monotonicity}. The
equality assumption and \eqref{eq:limit-Q} give $Q_k(0)=\lim_{s\to\infty}Q_k(s)$.
Therefore, for every $t\geq0$,
\begin{equation*}
    Q_k(0)\geq Q_k(t)\geq
    \lim_{s\to\infty}Q_k(s)=Q_k(0).
\end{equation*}
Hence $Q_k$ is constant and $D_tQ_k=0$ in distributions. Equation \eqref{eq:Q-defect} then implies
\begin{equation}\label{eq:equality-defects-zero}
    \mathfrak e_k=0,
    \qquad
    A_k=\frac{n-1+k}{n-1} F_k
    \quad\text{for almost every }t>0.
\end{equation}

We first show that there is no initial jump.  Assume that $  |\Om_0^+\setminus\Om|>0$.  Since $\partial\Om$ has zero $n$-dimensional measure, the open set $ \mathcal P:=\Om_0^+\setminus\overline\Om$ is nonempty.  The identity $\mathfrak e_k(\{0\})=0$ and the nonnegativity of both terms in \eqref{eq:Euclidean-defect} give
\begin{equation*}
    \partial_r\cdot\nu=1
    \quad\text{almost everywhere in }\mathcal P\setminus\{0\}.
\end{equation*}
Hence $\nu=\partial_r$ there.  Choose a ball $U$ compactly contained in $\mathcal P$ whose closure does not contain the origin.  Since $u$ is constant on $U$, equations \eqref{eq:calibration-div} and \eqref{eq:calibration} give
\begin{equation*}
    \diver\nu=0
    \quad\text{in }U.
\end{equation*}
On the other hand,
\begin{equation*}
    \diver \partial_r=\frac{n-1}{r}>0
    \quad\text{in }U,
\end{equation*}
which is a contradiction.  Thus $\Om_0^+=\Om$ up to a null set.

The second identity in \eqref{eq:equality-defects-zero} is equivalent to equality in \eqref{eq:divergence-estimate} for almost every $t$.  Lemma \ref{lem:divergence} shows that $\Om_t$ is a ball centered at the chosen origin for almost every $t$.  Choose such times $t_j\downarrow0$.  The absence of an initial jump and the right continuity of the weak flow give
\begin{equation*}
    \chi_{\Om_{t_j}}\longrightarrow\chi_\Om
    \quad\text{in }L^1_{\mathrm{loc}}.
\end{equation*}
If $R_j$ is the radius of $\Om_{t_j}$, then the area law \eqref{eq:area-growth} gives
\begin{equation*}
    R_j
    =\left(\frac{e^{t_j}|\Sig|}{|\Sn^{n-1}|}\right)^{\frac1{n-1}}
    \longrightarrow
    \left(\frac{|\Sig|}{|\Sn^{n-1}|}\right)^{\frac1{n-1}}.
\end{equation*}
Together with the local $L^1$ convergence, this shows that $\Om$ agrees almost everywhere with a ball centered at the chosen origin.  Since $\Om$ is a smooth domain, it is that centered ball.

Conversely, equality for a centered ball follows by direct computation.
\end{proof}

\subsection{Proof of the Euclidean Weinstock inequality}

We first record the consequence of the case $k=1$.

\begin{lemma}\label{lem:moment-lower}
Under the assumptions of Proposition \ref{prop:weighted-three-term},
\begin{equation*}
    \int_\Sig r^2\,d\mu
    \geq
    |B^n|^{-\frac2n}|\Sig|\,|\Om|^{\frac2n}.
\end{equation*}
Equality holds if and only if $\Om$ is a round ball centered at the origin.
\end{lemma}

\begin{proof}
The case $k=1$ of \eqref{eq:weighted-main} gives
\begin{equation}\label{eq:k1-three-term}
    \int_\Sig r\,d\mu
    \geq
    \frac{n-1}{n}|\Sn^{n-1}|^{-\frac1{n-1}}|\Sig|^{\frac n{n-1}}
    +|\Om|.
\end{equation}
By H\"older's inequality,
\begin{equation}\label{eq:holder-r}
    \int_\Sig r^2\,d\mu
    \geq
    \frac1{|\Sig|}\left(\int_\Sig r\,d\mu\right)^2.
\end{equation}
Using $|\Sn^{n-1}|=n|B^n|$, Young's inequality gives
\begin{equation*}
    |B^n|^{-1/n}|\Sig|\,|\Om|^{1/n}
    \leq
    |\Om|+\frac{n-1}{n}|\Sn^{n-1}|^{-\frac1{n-1}}|\Sig|^{\frac n{n-1}}.
\end{equation*}
Combining this with \eqref{eq:k1-three-term} gives
\begin{equation*}
    \int_\Sig r\,d\mu
    \geq
    |B^n|^{-1/n}|\Sig|\,|\Om|^{1/n}.
\end{equation*}
Equation \eqref{eq:holder-r} now yields the stated estimate.

Equality forces equality in Proposition \ref{prop:weighted-three-term}, H\"older's inequality, and Young's inequality.  Proposition \ref{prop:weighted-three-term} already implies that $\Om$ is a ball centered at the origin.  The converse is immediate.
\end{proof}


\begin{proof}[Proof of Theorem \ref{thm:weinstock}]
Choose the origin to be the boundary barycenter, so that
\begin{equation*}
    \int_\Sig x_i\,d\mu=0,
    \qquad i=1,\ldots,n.
\end{equation*}
Using $x_i$ as test functions in \eqref{eq:steklov-variational} gives
\begin{equation*}
    \sigma_1(\Om)\int_\Sig x_i^2\,d\mu
    \leq
    \int_\Om|\nabla x_i|^2\,dx
    =|\Om|.
\end{equation*}
Summing over $i$ gives
\begin{equation*}
    \sigma_1(\Om)\int_\Sig r^2\,d\mu
    \leq n|\Om|.
\end{equation*}
By Lemma \ref{lem:moment-lower},
\begin{equation*}
    \sigma_1(\Om)
    \leq
    n|B^n|^{\frac2n}|\Om|^{\frac{n-2}{n}}|\Sig|^{-1}.
\end{equation*}
Multiplying by $|\Sig|^{1/(n-1)}$ and using the isoperimetric inequality
\begin{equation*}
    |\Sig|\geq n|B^n|^{\frac1n}|\Om|^{\frac{n-1}{n}},
\end{equation*}
we obtain
\begin{equation*}
    \sigma_1(\Om)|\Sig|^{\frac1{n-1}}
    \leq
    (n|B^n|)^{\frac1{n-1}}
    =|\Sn^{n-1}|^{\frac1{n-1}}.
\end{equation*}
This is \eqref{eq:weinstock-main}, since $\sigma_1(B^n)=1$ for the unit ball.

If equality holds, then equality holds in Lemma \ref{lem:moment-lower}.  Therefore $\Om$ is a round ball.  Conversely, round balls attain equality.
\end{proof}


\begin{remark}\label{rem:Euclidean-endpoint}
Strict outward minimality is not required. Indeed, the minimizing-hull property and outward minimality give
\[
    P(\Omega_0^+)=P(\Omega)=|\Sigma|,
\]
so the area law remains $|\Sigma_t|=e^t|\Sigma|$. Moreover, \eqref{eq:Q-defect} gives
\begin{align*}
     (D_t Q_k)(\{0\})=&Q_k(0+)-Q_k(0)\\
     =&-|\Sigma|^{-\frac{n-1+k}{n-1}}\mathfrak e_k(\{0\}) \\
    =&-|\Sigma|^{-\frac{n-1+k}{n-1}} \left(
      \int_{\{u=0\}\setminus\overline\Omega}
        kr^{k-1}(1-\partial_r\cdot\nu)\,dx
      +\int_\Sigma r^k(1-\tau)\,d\mu
      \right) \leq 0.
\end{align*}
Here both integrals are nonnegative because $\partial_r\cdot\nu\leq1$ and $\tau\leq1$. Thus a possible enlargement at $t=0$ can only produce a downward jump of $Q_k$ and does not affect its monotonicity.
\end{remark}

\section{The hyperbolic case}\label{sec.4}

We now consider the hyperbolic case.  We first recall the radial functions $g$ and $h$ from Gu--Li--Wan \cite{GuLiWan2025}.  We then prove an endpoint weak-transport inequality for $\int_{\Sigma_t}g\dd\mu_t$.  To handle the denominator of the Gu--Li--Wan functional \eqref{s1.Mt}, we mollify the relevant quantities in time.  The weighted volume integral of a geodesic ball is convex as a function of the ball volume, so Jensen's inequality preserves the mass-transplantation estimate after mollification, including a possible contribution at $t=0$.  Finally, we insert the radial comparison estimates into the Steklov Rayleigh quotient.

\subsection{Notation from Gu--Li--Wan \cite{GuLiWan2025}}

We view the hyperbolic space as a warped product
\begin{equation*}
    (\HH^n,g_{\HH})=([0,\infty)\times\Sn^{n-1},dr^2+\lambda(r)^2g_{\Sn^{n-1}}),
    \qquad
    \lambda(r)=\sinh r.
\end{equation*}
After choosing an origin $O$, let $B_r$ and $S_r$ denote the geodesic ball and sphere of radius $r$ centered at the origin. We consider the volume-area ratio function 
\begin{align*}
    g(r)&:=\frac{|B_r|}{|S_r|}=\frac{1}{\lambda(r)^{n-1}}\int_0^r\lambda(s)^{n-1}\,ds,
 %   a(r)&:=\frac{\lambda'(r)g(r)}{\lambda(r)}.
\end{align*}
Then the first positive Steklov eigenvalue of the geodesic ball $B_r\subset \mathbb{H}^n$ satisfies 
\begin{equation*}
    \sigma_1(B_r)=\frac{g'(r)}{g(r)}. 
\end{equation*}
Denote 
\begin{equation}\label{s4.q}
    q(r):=(g'(r))^2+(n-1)\frac{g(r)^2}{\lambda(r)^2}.
\end{equation}
The identities used below are
\begin{equation}\label{s4.q.div}
    g'=1-(n-1)\frac{\lambda'g}{\lambda},
    \qquad
    \diver(g\partial_r)=1,
    \qquad
    \diver(gg'\partial_r)=q.
\end{equation}
The analytic properties needed below are collected in
\cite[Propositions~3.4 and~3.5]{GuLiWan2025}: $g$ is nonnegative,
strictly increasing and concave, with
\begin{equation*}
    g(0)=0,
    \qquad g'(0)=\frac1n,
    \qquad \lim_{r\to\infty}g(r)=\frac1{n-1},
    \qquad \lim_{r\to\infty}g'(r)=0.
\end{equation*}
In particular, concavity makes $g'$ non-increasing and
$0<g'(r)\leq1/n$ for finite $r$.  The first identity in
\eqref{s4.q.div} also gives
\begin{equation*}
    \frac{\lambda'g}{\lambda}=\frac{1-g'}{n-1},
\end{equation*}
so this ratio is non-decreasing and
\begin{equation}\label{eq:hyp-a-range}
    \frac1n\leq \frac{\lambda'(r)g(r)}{\lambda(r)} \leq\frac1{n-1}.
\end{equation}
The radial function $x\mapsto g(r(x))$ is locally Lipschitz near the origin and $ \nabla g=g'(r)\partial_r$ 
almost everywhere.

For a finite-perimeter set $E\subset\HH^n$ define
\begin{align*}
    J(E):=\int_E\frac{\lambda'g}{\lambda}\,dv.
\end{align*}
 Set
\begin{equation*}
    J_B(r):=J(B_r),
    \qquad
    V_B(r):=|B_r|.
\end{equation*}
Following Gu--Li--Wan \cite[\S 4.1]{GuLiWan2025}, define $h:(0,\infty)\to(0,\infty)$ by
\begin{equation*}
    h(J_B(r))=\frac1{|S_r|}.
\end{equation*}
Indeed,
\begin{equation*}
    J_B'(r)
    =\frac{\lambda'(r)g(r)}{\lambda(r)}|S_r|>0,
\end{equation*}
and \eqref{eq:hyp-a-range} gives
\begin{equation*}
    J_B(r)\geq\frac1nV_B(r)\longrightarrow\infty
    \qquad\text{as }r\to\infty.
\end{equation*}
Hence $J_B$ is a smooth increasing bijection from $(0,\infty)$ onto
$(0,\infty)$. Consequently,
\begin{equation*}
    h(s)=\frac{1}{\left|S_{J_B^{-1}(s)}\right|}
\end{equation*}
is smooth on $(0,\infty)$. Lemma~4.1 of \cite{GuLiWan2025} states that $h$ is positive, decreasing and log-convex, and gives
\begin{equation}\label{eq:hyp-h-log-derivative}
    \frac{h'(J_B(r))}{h(J_B(r))}
    =-\frac{n-1}{V_B(r)}.
\end{equation}
For later use, log-convexity implies
\begin{equation}\label{eq:hyp-hhprime-monotone}
    \bigl(hh'\bigr)'=(h')^2+hh''\geq0,
\end{equation}
so $s\mapsto h(s)h'(s)$ is non-decreasing on $(0,\infty)$. Since $\lambda'g/\lambda$ is increasing, mass transplantation (see \cite[\S 4]{FL21} and \cite{We56}) gives
\begin{equation}\label{eq:hyp-J-mass-transplant}
    J(E)\geq J(B_r),
    \qquad \text{when}~~ |B_r|=|E|.
\end{equation}
The function $J_B$, viewed as a function of ball volume, is convex because
\begin{equation*}
    \frac{dJ_B}{dV_B}=\frac{\lambda'(r)g(r)}{\lambda(r)}
\end{equation*}
and $\lambda'g/\lambda$ is increasing.

\subsection{A weak transport inequality for \texorpdfstring{$\int_{\Sigma_t}g\dd\mu_t$}{int Sigma g}}

Let $\Omega\subset\HH^n$ be a bounded outward-minimizing domain with smooth boundary $\Sigma$.  Let $\{\Omega_t\}_{t\geq0}$ be the proper weak inverse mean curvature flow starting from $\Omega$.  We set $\Omega_0=\Omega$ and $\Sigma_0=\Sigma$.  The area growth law is
\begin{equation*}
    |\Sigma_t|=e^t|\Sigma|,
    \qquad t\geq0.
\end{equation*}

With respect to the chosen origin $O$, define, for $t>0$,
\begin{equation*}
    V(t):=|\Omega_t|,\qquad
    J(t):=\int_{\Omega_t}\frac{\lambda'g}{\lambda}\,dv,\qquad
    C(t):=\int_{\Omega_t}g'\,dv,
\end{equation*}
and, for almost every $t>0$, define
\begin{equation*}
    G(t):=\int_{\Sigma_t}g\,d\mu_t.
\end{equation*}
At $t=0$, prescribe
\begin{equation*}
    G(0):=\int_\Sigma g\,d\mu,\qquad
    V(0):=|\Omega|,\qquad
    J(0):=\int_\Omega\frac{\lambda'g}{\lambda}\,dv,\qquad
    C(0):=\int_\Omega g'\,dv.
\end{equation*}
The identity \eqref{s4.q.div} gives
\begin{equation*}
    C(t)=V(t)-(n-1)J(t),
    \qquad t\geq0.
\end{equation*}
Since the weak-flow sets $\Omega_t$ are nested and both
$\lambda'g/\lambda$ and $g'$ are nonnegative, the functions
$V$, $J$ and $C$ are non-decreasing on $[0,\infty)$. Whenever
distributional derivatives or time convolutions are used below, all four functions are extended to $t<0$ by their prescribed values at $t=0$.

\begin{lemma}\label{lem:hyp-weak-transport-G}
The function $G$ belongs to $L^1_{\mathrm{loc}}([0,\infty))$.  After the constant extensions described above, the following inequality holds in distributions on $\mathbb R$:
\begin{equation}\label{eq:hyp-measure-transport}
    D_tG\leq G\,dt+D_tC.
\end{equation}
Equivalently, for every nonnegative $\eta\in C_c^1([0,\infty))$,
\begin{equation}\label{eq:hyp-transport-test}
    -\int_0^\infty\eta'(t)G(t)\,dt-\eta(0)G(0)
    \leq
    \int_0^\infty\eta(t)G(t)\,dt
    -\int_0^\infty\eta'(t)C(t)\,dt-\eta(0)C(0).
\end{equation}
\end{lemma}


\begin{thebibliography}{99}
\bibitem[AFM2022]{AFM2022} V. Agostiniani, M. Fogagnolo, and L. Mazzieri, \emph{Minkowski inequalities via nonlinear potential theory}, Arch. Ration. Mech. Anal. \textbf{244} (2022), 51--85.
\bibitem[BFNT2021]{BFNT2021} D. Bucur, V. Ferone, C. Nitsch, and C. Trombetti, \emph{Weinstock inequality in higher dimensions}, J. Differential Geom. \textbf{118} (2021), 1--21.
\bibitem[CGGS24]{CGGS24} B. Colbois, A. Girouard, C. Gordon, and D. Sher, \emph{Some recent developments on the Steklov eigenvalue problem}, Rev. Mat. Complut. \textbf{37} (2024), 1--161.
\bibitem[ChenFrid1999]{ChenFrid1999} G.-Q. Chen and H. Frid, \emph{Divergence-measure fields and hyperbolic conservation laws}, Arch. Ration. Mech. Anal. \textbf{147} (1999), 89--118.
\bibitem[FL21]{FL21} P. Freitas and R.~S. Laugesen, \emph{From Neumann to Steklov and beyond, via Robin: the Weinberger way}, Amer. J. Math. {\bf 143} (2021), no.~3, 969--994
\bibitem[FM22]{FM22} M. Fogagnolo and L. Mazzieri, \emph{Minimising hulls, $p$-capacity and isoperimetric inequality on complete Riemannian manifolds}, J. Funct. Anal. \textbf{283} (2022), Paper No. 109638.
\bibitem[FS2014]{FS2014} A. Freire and F. Schwartz, \emph{Mass-capacity inequalities for conformally flat manifolds with boundary}, Comm. Partial Differential Equations \textbf{39} (2014), 98--119.
\bibitem[FS2019]{FS2019} A. Fraser and R. Schoen, \emph{Shape optimization for the Steklov problem in higher dimensions}, Adv. Math. \textbf{348} (2019), 146--162.
\bibitem[Fra20]{Fra20} A. Fraser, \emph{Extremal eigenvalue problems and free boundary minimal surfaces in the ball}, in \emph{Geometric analysis}, Lecture Notes in Math., vol. 2263, Springer, Cham, 2020, 1--40.
\bibitem[GLW26]{GLW26} P. Gu, H. Li, and Y. Wan, \emph{Weinstock inequality in hyperbolic space II}, preprint submitted to Proceedings of ICCM.
\bibitem[GP2017]{GP2017} A. Girouard and I. Polterovich, \emph{Spectral geometry of the Steklov problem}, J. Spectr. Theory \textbf{7} (2017), 321--359.
\bibitem[GR2020]{GR2020} F. Gir\~ao and D. Rodrigues, \emph{Weighted geometric inequalities for hypersurfaces in sub-static manifolds}, Bull. Lond. Math. Soc. \textbf{52} (2020), 121--136.
\bibitem[Gerhardt1990]{Gerhardt1990} C. Gerhardt, \emph{Flow of nonconvex hypersurfaces into spheres}, J. Differential Geom. \textbf{32} (1990), 299--314.
\bibitem[Gerhardt2011]{Gerhardt2011} C. Gerhardt, \emph{Inverse curvature flows in hyperbolic space}, J. Differential Geom. \textbf{89} (2011), 487--527.
\bibitem[GornyMazon2024]{GornyMazon2024} W.~G{\'o}rny and J.~M. Maz{\'o}n, \emph{The Anzellotti--Gauss--Green formula and least gradient functions in metric measure spaces}, Commun. Contemp. Math. \textbf{26} (2024), no.~6, Paper No.~2350027, 42 pp.
\bibitem[GuLiWan2025]{GuLiWan2025} P. Gu, H. Li, and Y. Wan, \emph{Weinstock inequality in hyperbolic space}, J. Funct. Anal. \textbf{289} (2025), 111155.
\bibitem[GuThesis2026]{GuThesis2026} P. Gu, \emph{Isoperimetric problems and optimal domains for Steklov eigenvalues in space forms}, Ph.D. thesis, Tsinghua University, 2026.
\bibitem[HI2001]{HI2001} G. Huisken and T. Ilmanen, \emph{The inverse mean curvature flow and the Riemannian Penrose inequality}, J. Differential Geom. \textbf{59} (2001), 353--437.
\bibitem[HI2008]{HI2008} G. Huisken and T. Ilmanen, \emph{Higher regularity of the inverse mean curvature flow}, J. Differential Geom. \textbf{80} (2008), 433--451.
\bibitem[Harvie2026]{Harvie2026} B. Harvie, \emph{On weak inverse mean curvature flow and Minkowski-type inequalities in hyperbolic space}, Calc. Var. Partial Differential Equations \textbf{65} (2026), Paper No. 212. %\href{https://doi.org/10.1007/s00526-026-03384-4}{doi:10.1007/s00526-026-03384-4}.
\bibitem[Huisken2009Lecture]{Huisken2009Lecture} G. Huisken, \emph{Marston Morse: Inverse mean curvature flow and isoperimetric inequalities}, video lecture, Institute for Advanced Study, March 20, 2009, \href{https://www.ias.edu/video/marston-morse-inverse-mean-curvature-flow-and-isoperimetric-inequalities}{IAS video archive}.
\bibitem[KwongWei2023]{KwongWei2023} K.-K. Kwong and Y. Wei, \emph{Geometric inequalities involving three quantities in warped product manifolds}, Adv. Math. \textbf{430} (2023), Paper No. 109213.
\bibitem[LMP2023]{LMP2023} M. Levitin, D. Mangoubi, and I. Polterovich, \emph{Topics in spectral geometry}, Graduate Studies in Mathematics, vol. 237, American Mathematical Society, Providence, RI, 2023.
\bibitem[Maggi2012]{Maggi2012} F. Maggi, \emph{Sets of finite perimeter and geometric variational problems}, Cambridge Studies in Advanced Mathematics, vol. 135, Cambridge University Press, 2012.
\bibitem[Urbas1990]{Urbas1990} J. I. E. Urbas, \emph{On the expansion of starshaped hypersurfaces by symmetric functions of their principal curvatures}, Math. Z. \textbf{205} (1990), 355--372.
\bibitem[W1954]{W1954} R. Weinstock, \emph{Inequalities for a classical eigenvalue problem}, J. Rational Mech. Anal. \textbf{3} (1954), 745--753.
\bibitem[We56]{We56} H.~F. Weinberger, \emph{An isoperimetric inequality for the $N$-dimensional free membrane problem}, J. Rational Mech. Anal. {\bf 5} (1956), 633--636
\end{thebibliography}
\end{document}
